\section{Mathematical setup}
\label{sec:setup}

\subsection{Reference process and learned drift}

Let $\HH$ be a separable real Hilbert space.  Let
$\mathcal C:\HH\to\HH$ be positive, self-adjoint, injective, and
trace-class.  Write
\[
  \mathcal C\phi_j=\lambda_j\phi_j,\qquad
  \lambda_1\geq\lambda_2\geq\cdots>0,
  \qquad \sum_{j\geq1}\lambda_j<\infty ,
\]
for an orthonormal eigenbasis, and let
$\mu_{\mathrm{ref}}=\mathcal N(0,\mathcal C)$.  Its
Cameron--Martin space is
\[
  H_{\mathcal C}=\operatorname{Ran}(\mathcal C^{1/2}),\qquad
  \norm{h}_{H_{\mathcal C}}
  =\norm{\mathcal C^{-1/2}h}_{\HH}.
\]
The inverse in this expression is the inverse on
$\operatorname{Ran}(\mathcal C^{1/2})$.

On a filtered probability space
$(\Omega,\mathcal F,(\mathcal F_t)_{t\leq T},\mathbb Q)$, let $W$ be a
cylindrical Wiener process on $\HH$.  The stationary reference process
solves
\begin{equation}
  \mathrm dX_t=-\tfrac12 X_t\,\mathrm dt
      +\mathcal C^{1/2}\,\mathrm dW_t,
  \qquad X_0\sim\mu_{\mathrm{ref}},
  \label{eq:reference-ou}
\end{equation}
with $X_0$ independent of $W$.  Thus
$\mathcal L_{\mathbb Q}(X_t)=\mu_{\mathrm{ref}}$ for every $t$.

\begin{assumption}[Cameron--Martin drift]
\label{ass:drift}
The map $\drift:[0,T]\times\HH\to H_{\mathcal C}$ is jointly
measurable.  There are constants $L,M<\infty$ such that, for every
$t\in[0,T]$ and $u,v\in\HH$,
\begin{align}
  \norm{\drift(t,u)-\drift(t,v)}_{H_{\mathcal C}}
    &\leq L\norm{u-v}_{\HH}, \label{eq:drift-lipschitz}\\
  \norm{\drift(t,u)}_{H_{\mathcal C}}
    &\leq M(1+\norm{u}_{\HH}). \label{eq:drift-growth}
\end{align}
\end{assumption}

The perturbed process is
\begin{equation}
  \mathrm dU_t
  =\bigl[-\tfrac12 U_t+\drift(t,U_t)\bigr]\,\mathrm dt
   +\mathcal C^{1/2}\,\mathrm dW_t,
  \qquad U_0\sim\mu_{\mathrm{ref}},
  \label{eq:prior-sde}
\end{equation}
and the prior is its terminal law
\begin{equation}
  \muz=\mathcal L(U_T).
  \label{eq:prior-terminal-law}
\end{equation}
\Cref{ass:drift} implies existence, pathwise uniqueness, and finite
moments of every fixed order used below; see
\Cref{lem:sde-moments}.

Set
\[
  a(t,u)=\mathcal C^{-1/2}\drift(t,u).
\]
The following condition is deliberately stated on the reference path.

\begin{assumption}[Finite-horizon change of measure]
\label{ass:girsanov}
The stochastic exponential
\begin{equation}
  \mathcal E_T
  =
  \exp\!\left\{
    \int_0^T\inner{a(t,X_t)}{\mathrm dW_t}_{\HH}
    -\frac12\int_0^T\norm{a(t,X_t)}_{\HH}^2\,\mathrm dt
  \right\}
  \label{eq:girsanov-exponential}
\end{equation}
is a true martingale and
$\mathbb E_{\mathbb Q}\mathcal E_T^2<\infty$.
\end{assumption}

Under this assumption, the law of $X$ after the change of measure
$\mathrm d\mathbb Q^\theta=\mathcal E_T\,\mathrm d\mathbb Q$ is the
law of \eqref{eq:prior-sde}.  In particular,
\begin{equation}
  \muz\sim\mu_{\mathrm{ref}},\qquad
  \rho(x):=\frac{\mathrm d\muz}{\mathrm d\mu_{\mathrm{ref}}}(x)
  =\mathbb E_{\mathbb Q}
      [\mathcal E_T\mid X_T=x],
  \label{eq:terminal-density}
\end{equation}
for a suitable version of the conditional expectation, and
$\rho\in L^2(\mu_{\mathrm{ref}})$; see
\Cref{lem:girsanov-terminal}.
Bounded $a$ is a simple sufficient condition for the martingale part of
\Cref{ass:girsanov}.  Affine growth alone is not a substitute for
Novikov's or another valid martingale criterion, so the condition is
kept explicit.

\begin{remark}[Relation to a reverse-time score]
\label{rem:not-reverse-score}
Equation \eqref{eq:prior-sde} defines a learned-drift prior.  It is not,
by itself, the reverse-time SDE associated with
\eqref{eq:reference-ou}.  In the reverse-time formulation, the time
orientation and drift contain the score of the relevant noised
marginal, characterized through conditional expectations
\citep{pidstrigach2023}.  A network trained by denoising score matching
may be used to parameterize $\drift$, but identifying the trained
network with that reverse-time score requires a separate approximation
theorem.  No such identification is used here.
\end{remark}

\subsection{Observation model and posterior}

Let $\G:\HH\to\R^k$ be bounded and linear, let
$\Gamobs\in\R^{k\times k}$ be positive definite, and fix
$\sigobs>0$.  The data satisfy
\begin{equation}
  y=\G\udag+\sigobs\noise,\qquad
  \noise\sim\mathcal N(0,\Gamobs).
  \label{eq:forward}
\end{equation}
The negative log-likelihood, up to an additive constant independent of
$u$, is
\begin{equation}
  \Phim_{\sigobs}(u;y)
  =\frac{1}{2\sigobs^2}
    \norm{\Gamobs^{-1/2}(y-\G u)}_{\R^k}^2 .
  \label{eq:potential}
\end{equation}
Bayes' formula defines
\begin{equation}
  \frac{\mathrm d\muy}{\mathrm d\muz}(u)
  =\frac{1}{Z(y)}
     \exp\{-\Phim_{\sigobs}(u;y)\},\qquad
  Z(y)=\int_{\HH}\exp\{-\Phim_{\sigobs}(u;y)\}\,\muz(\mathrm du).
  \label{eq:posterior}
\end{equation}
The noise factor $\sigobs^{-2}$ in \eqref{eq:potential} is retained
throughout; it is essential when studying small-noise limits.

For probability measures $\mu,\nu$ dominated by $\xi$, we use
\begin{equation}
  \dH^2(\mu,\nu)
  =\frac12\int
    \left(
      \sqrt{\frac{\mathrm d\mu}{\mathrm d\xi}}
      -\sqrt{\frac{\mathrm d\nu}{\mathrm d\xi}}
    \right)^2\,\mathrm d\xi .
  \label{eq:hellinger}
\end{equation}
The value does not depend on the choice of $\xi$.

\subsection{Spectral projections}

Let
\[
  P_Nu=\sum_{j=1}^N\inner{u}{\phi_j}\phi_j,\qquad
  \HN=P_N\HH,
\]
and define the covariance-tail scale
\begin{equation}
  r_N^2
  =\operatorname{Tr}\bigl((I-P_N)\mathcal C\bigr)
  =\sum_{j>N}\lambda_j .
  \label{eq:tail-scale}
\end{equation}
Because $\mathcal C$ is trace-class, $r_N\to0$.  The approximating
prior used below remains a measure on all of $\HH$; it does not replace
the unresolved coordinates by zero.
